\section{Relation to earlier work and the next physical test}
\label{sec:completion-positioning}

The contribution of access completion is a chain of usable certificates:
identify an inaccessible target, specify the additional physical operation,
bound what its data can establish, and expose the assumptions under which the
conclusion fails. Several links in this chain are established mathematics.
Separating those foundations from the results derived here makes the proposed
physical test easier to assess and reproduce.

\paragraph{Established foundations.}
Conservation-law sensor placement, spanning-tree constructions, cost, and
precision constraints have a substantial earlier
literature~\cite{ac_bagajewicz,ac_he,ac_xu}. The cotree and target-cut results
are explicit network-flow specializations. Perturbation design based on
graph-theoretic identifiability also predates this work~\cite{ac_gross}.
Target-specific allocation belongs to optimal experimental
design~\cite{ac_elfving,ac_harman}; adding resource constraints to terminal cuts
has established complexity consequences~\cite{ac_aissi}. The particular
calibration-floor objective studied here therefore requires its stated error
model and direct proof. It is not a universal thermodynamic energy law.
Likewise, lumpability and approximate trace guarantees underpin the classical
interface construction~\cite{ac_buchholz,ac_bian}. Quantum equivalence under
composition has its own established theory~\cite{ac_lalire,ac_deng}.

\paragraph{The precise response result.}
Stationary barrier-response identities are
established~\cite{ac_aslyamov_response,ac_aslyamov_static}.
Finite-amplitude stationary formulas already include parallel physical
channels~\cite{ac_bao_nonlinear}; transient response relations are also
available~\cite{ac_zheng}. Monotonicity of a stationary response as a barrier
strength varies~\cite{ac_floyd} concerns a different variable from the pulse
duration used here. First-passage response theory supplies further close
connections~\cite{ac_liang_budget,ac_bao_firstpassage}.

The candidate extension is the following all-time statement. Starting from
native stationarity, suppress one bidirectional physical channel by a fixed
factor $r<1$ while holding the other rates fixed. For a singleton endpoint
detector, the response factors as
\begin{equation}
 \Delta p_j(t)=-(1-r)J_e H_{ij}^{r}(t),
 \qquad H_{ij}^{r}(t)>0\quad(t>0).
\end{equation}
The first-hit proof uses the strong Markov property and ordinary occupation
potentials, within established Markov potential theory~\cite{ac_coolen}.
Its useful consequence is a current-sign certificate at every positive pulse
duration, with explicit bounds using only endpoint exit-rate ceilings.
It does not require reconstructing the hidden network. A protected coarse
detector also qualifies when a separately justified entrance-gateway condition
holds; arbitrary aggregation does not preserve the conclusion.

The driven-qubit calculation gives an exact parameter boundary at which this
classical response restriction fails. Connections between coherent dynamics
and failure of classical stochastic descriptions are already well
developed~\cite{ac_smirne,ac_prech}. Here the distinguishing feature is the
specified intervention and endpoint-response inequality. The result rejects
the calibrated classical model class; its assumptions cannot be inferred
from two response averages alone.

\paragraph{Priority and evidence limits.}
The proofs and computational certificates establish the stated mathematical
claims within their models. They establish neither historical priority nor
laboratory performance. The particularly close sensor-design article by
Prakash~\cite{ac_prakash} was available only through primary metadata and
indexed excerpts in this audit. The higher-cumulant response article of Bao
and Liang~\cite{ac_bao_higher} was inspected at abstract level only. Their full
texts remain necessary overlap checks. The search of older finite-time
potential-kernel results is also incomplete. A detailed source map records
which full texts, sections, or abstracts were inspected. These physical-response examples have no new empirical validation or external
peer-review decision. Chapter~\ref{ch:ci-empirical} separately reports an
exploratory reanalysis of published optical-tweezer measurements.

\paragraph{The next validation gate.}
The highest-leverage physical test is a controlled two-level platform with
independently adjustable incoherent excitation and decay channels. It should
first establish stationary preparation, population-readout calibration and
leakage bounds; then independently verify that the pulse rescales both rates
of the selected channel while leaving the Hamiltonian and other channels
within declared tolerances. This calibration is part of the experiment,
not an assumption justified by fitting the desired sign reversal.

The measurement then compares native and pulsed endpoint populations at two
prespecified durations, with randomized acquisition order and independent
repreparation. The incoherent control should obey the classical sign
restriction. Increasing the calibrated coherent-drive ratio across the
derived threshold should produce the predicted reversed response. Sampling,
preparation, drift, detector, and intervention error budgets must be small
enough that the two response intervals exclude zero with opposite signs.
Repeating the experiment with deliberately nonselective control or coarse
readout tests the documented failure modes rather than merely reproducing
the favorable case.

Passing this gate would validate a practical way to test a structural model
assumption using a small set of interventions. It would also determine which
parts of the calibration dominate sensitivity. A failure would locate a
specific mismatch between the physical platform and the model. Either
outcome is more informative than projecting the framework directly onto
unresolved problems of large-scale computing or biological flux.
